\paragrafos{Προβλήματα με ποσοστά: εκπτώσεις, φόροι, τόκοι}{Α.5.2, «Προβλήματα με ποσοστά»}

\begin{anagke}
Ιανουάριος και στα μαγαζιά της γειτονιάς κρέμονται πινακίδες με εκπτώσεις. Ο Νίκος θέλει αθλητικά των $90$ €. Η πινακίδα γράφει «Έκπτωση $30\%$», και από κάτω «Επιπλέον $20\%$ στο ταμείο». Ο Βασίλης κάνει το λογαριασμό με σιγουριά: «Τριάντα και είκοσι κάνουν πενήντα. Μισή τιμή! Σαράντα πέντε ευρώ». Ο Νίκος χαίρεται, αλλά η Πολυξένη, που έχει μάθει να μην εμπιστεύεται λογαριασμούς πριν τους ελέγξει, γράφει: «Πρώτα το $30\%$ του $90$ είναι $27$, άρα η τιμή γίνεται $63$. Μετά το $20\%$ είναι του $63$ και όχι του $90$: $12{,}60$ €. Η τιμή γίνεται $50{,}40$».

«Πενήντα ευρώ και σαράντα λεπτά. Όχι σαράντα πέντε», μουρμουρίζει ο Νίκος, «πέντε ευρώ παραπάνω;». Η Αγγελική συμπληρώνει: «Δεν σε κοροϊδεύουν. Το $20\%$ υπολογίζεται πάνω στη νέα τιμή».

Αυτά τα προβλήματα, ο φόρος της απόδειξης, η έκπτωση στο ταμείο, ο τόκος της τράπεζας, και η ποσοστιαία αύξηση ή μείωση μιας τιμής, είναι η αιτία που ο καθένας πρέπει να ξέρει ποσοστά. Και το ερώτημα «πόσο πρέπει να αυξηθεί μια τιμή που έπεσε $20\%$ για να γυρίσει πίσω;» δεν έχει την απάντηση που περιμένει κανείς.
\end{anagke}

\begin{lexi}{έκπτωση, τόκος, επιτόκιο, ΦΠΑ}{έκπτωση, τόκος, ΦΠΑ}
Η \lexiword{έκπτωση} είναι το «εκ-πίπτω», δηλαδή ό,τι «πέφτει» από την τιμή. Ο \lexiword{τόκος} είναι τα χρήματα που «γεννιούνται» από ένα κεφάλαιο που έχουμε δώσει στην τράπεζα (τόκος από το «τίκτω»: γεννάω). Το \lexiword{επιτόκιο} είναι το ποσοστό του τόκου. Το \lexiword{κεφάλαιο} (από το «κεφαλή») είναι το βασικό ποσό. Το \lexiword{ΦΠΑ} είναι ο Φόρος Προστιθέμενης Αξίας, που προστίθεται στην τιμή ενός προϊόντος.
\end{lexi}

\begin{metafrasi}
Τιμή μετά από έκπτωση $30\%$. & $\text{τιμή}\cdot\left(1-\dfrac{30}{100}\right)=0{,}7\cdot\text{τιμή}$\\
Τιμή μετά από αύξηση $15\%$. & $\text{τιμή}\cdot1{,}15$\\
Ποσοστιαία μεταβολή. & $\dfrac{\text{νέα}-\text{παλιά}}{\text{παλιά}}$ σε μορφή ποσοστού\\
Τόκος ενός έτους. & $\text{κεφάλαιο}\cdot\dfrac{\text{επιτόκιο}}{100}$\\
Ο ΦΠΑ μέσα στην τιμή με ΦΠΑ $24\%$. & $\dfrac{24}{124}$ της τελικής τιμής\\
\end{metafrasi}

\tmima{Η θεωρία}

\tmimasmall{Έκπτωση και αύξηση}

Για έκπτωση $\alpha\%$ ή αύξηση $\alpha\%$ σε μια τιμή:
$$\text{τιμή μετά την έκπτωση}=\text{τιμή}-\dfrac{\alpha}{100}\cdot\text{τιμή}=\left(1-\dfrac{\alpha}{100}\right)\cdot\text{τιμή},$$
$$\text{τιμή μετά την αύξηση}=\left(1+\dfrac{\alpha}{100}\right)\cdot\text{τιμή}.$$
Αν η έκπτωση είναι $20\%$, πολλαπλασιάζω την τιμή με $0{,}8$. Αν η αύξηση είναι $15\%$, με $1{,}15$.

\tmimasmall{Ποσοστιαία μεταβολή}

Για να βρω πόσο \% ανέβηκε ή έπεσε μια τιμή από $\pi$ σε $\nu$ βρίσκω τη διαφορά και τη συγκρίνω με την \emph{αρχική} τιμή: $\dfrac{\nu-\pi}{\pi}$. Από $40$ σε $52$: $\frac{12}{40}=0{,}3=30\%$ αύξηση. Από $52$ σε $39$: $\frac{13}{52}=0{,}25=25\%$ μείωση.

\tmimasmall{Βρίσκω την αρχική τιμή}

Όταν ξέρω τη νέα τιμή και το ποσοστό, γράφω εξίσωση: μετά από έκπτωση $25\%$ η τιμή είναι $45$ €. Αν $x$ η αρχική, $0{,}75x=45$, άρα $x=45:0{,}75=60$ €. Και, γενικά, αν το $\alpha\%$ ενός αριθμού είναι $\mu$, ο αριθμός είναι $\mu:\frac{\alpha}{100}$.

\tmimasmall{Διαδοχικές μεταβολές}

Δύο μεταβολές δεν προστίθενται: κάθε νέο ποσοστό υπολογίζεται πάνω στη νέα τιμή. Έκπτωση $30\%$ και μετά $20\%$ σημαίνει πολλαπλασιασμός με $0{,}7\cdot0{,}8=0{,}56$, δηλαδή τελική έκπτωση $44\%$ και όχι $50\%$.

\tmimasmall{ΦΠΑ μέσα στην τιμή}

Αν μια τιμή περιέχει ΦΠΑ $24\%$, η τιμή χωρίς ΦΠΑ είναι $100$ και η τελική $124$ (ανά $100$). Άρα ο ΦΠΑ είναι τα $\dfrac{24}{124}$ της τελικής τιμής και όχι το $24\%$ της. Για τελική τιμή $372$ €, ο ΦΠΑ είναι $372\cdot\frac{24}{124}=72$ € και η αξία χωρίς ΦΠΑ είναι $372-72=300$ €.

\tmimasmall{Τόκος}

Ο \textbf{τόκος} ενός κεφαλαίου για ένα έτος είναι $\text{κεφάλαιο}\cdot\text{επιτόκιο}$. Για κλάσμα του έτους παίρνουμε το αντίστοιχο κλάσμα: για $6$ μήνες τα $\frac{6}{12}$. Αν οι τόκοι προστίθενται στο κεφάλαιο και στο τέλος κάθε έτους (\emph{κεφαλαιοποίηση}), το επόμενο έτος ο τόκος υπολογίζεται στο νέο κεφάλαιο.

\tmimasmall{Παράσταση ποσοστών με διαγράμματα}

Τα ποσοστά ενός συνόλου σχεδιάζονται με \textbf{κυκλικό διάγραμμα}, όπου ολόκληρος ο κύκλος είναι το $100\%$ και κάθε τομέας έχει εμβαδόν ανάλογο του ποσοστού (το $1\%$ αντιστοιχεί σε γωνία $3{,}6^\circ$), ή με \textbf{ραβδόγραμμα}, όπου το ύψος κάθε ράβδου είναι το ποσοστό. Τα ποσοστά ενός πλήρους συνόλου έχουν πάντα άθροισμα $100\%$.

\begin{center}
\begin{tikzpicture}[scale=0.9]
\def\r{1.7}
\colorlet{c1}{gblue}\colorlet{c2}{gblue!70}\colorlet{c3}{gblue!45}\colorlet{c4}{gblue!22}
\fill[c1] (0,0)--(90:\r) arc(90:-54:\r)--cycle;
\fill[c2] (0,0)--(-54:\r) arc(-54:-180:\r)--cycle;
\fill[c3] (0,0)--(-180:\r) arc(-180:-234:\r)--cycle;
\fill[c4] (0,0)--(-234:\r) arc(-234:-270:\r)--cycle;
\draw[white,thick] (0,0)--(90:\r) (0,0)--(-54:\r) (0,0)--(-180:\r) (0,0)--(-234:\r);
\draw (0,0) circle (\r);
\node at (18:1.05) {\small$40\%$};
\node at (-117:1.0) {\small$35\%$};
\node at (-207:1.05) {\small$15\%$};
\node[font=\scriptsize] at (-252:1.15) {$10\%$};
\begin{scope}[xshift=5.6cm,yshift=-1.4cm]
\foreach \y/\p/\c/\t in {0/40/c1/μεταφορά,0.75/35/c2/διαμονή,1.5/15/c3/φαγητό,2.25/10/c4/ψυχαγωγία}{
\fill[\c] (0,\y)rectangle(\p*0.07,\y+0.5);
\node[right,font=\small] at (\p*0.07+0.1,\y+0.25) {$\p\%$ \t};}
\end{scope}
\end{tikzpicture}
\end{center}

Οι δαπάνες των $600$ € της εκδρομής σε κυκλικό διάγραμμα και σε ραβδόγραμμα: μεταφορά $40\%$ ($240$ €), διαμονή $35\%$ ($210$ €), φαγητό $15\%$ ($90$ €), ψυχαγωγία $10\%$ ($60$ €).

\begin{prosoxi}
\begin{itemize}[leftmargin=1.2em,itemsep=0.3ex,topsep=2pt]
\item Εκπτώσεις και αυξήσεις δεν προστίθενται: $30\%$ και μετά $20\%$ δεν κάνουν $50\%$, αλλά $44\%$.
\item Το ποσοστό μεταβολής υπολογίζεται πάντα επί της \emph{αρχικής} τιμής. Μια τιμή που ανέβηκε από $100$ σε $125$ ($+25\%$) πρέπει να πέσει κατά $20\%$ (όχι $25\%$) για να γυρίσει στα $100$.
\item Ο ΦΠΑ μέσα στην τελική τιμή δεν είναι το $24\%$ της τελικής τιμής, αλλά τα $\frac{24}{124}$.
\item Στο κυκλικό διάγραμμα τα ποσοστά πρέπει να κάνουν σύνολο $100\%$.
\end{itemize}
\end{prosoxi}

\tmima{Λυμένα παραδείγματα}

\paradeigma{Οι διαδοχικές εκπτώσεις}
\tixero Αθλητικά $90$ €, έκπτωση $30\%$ και μετά $20\%$.\par
\tithelo Η τελική τιμή και η συνολική έκπτωση.\par
\possindeo Μετά την πρώτη: $0{,}7\cdot90=63$ €. Μετά τη δεύτερη: $0{,}8\cdot63=50{,}4$ €. Συνολική έκπτωση: $90-50{,}4=39{,}6$ €, δηλαδή $\frac{39{,}6}{90}=0{,}44=44\%$.\par
\apantisi $50{,}40$ € και συνολική έκπτωση $44\%$.\par
\elegxos $0{,}7\cdot0{,}8=0{,}56$ και $0{,}56\cdot90=50{,}4$. Το $50\%$ θα έδινε $45$ €.

\paradeigma{Ο ΦΠΑ μέσα στην τιμή}
\tixero Ένα ποδήλατο κοστίζει $372$ € με ΦΠΑ $24\%$.\par
\tithelo Ο ΦΠΑ και η αξία χωρίς ΦΠΑ.\par
\possindeo Αν $x$ η αξία χωρίς ΦΠΑ, $1{,}24x=372$, άρα $x=372:1{,}24=300$. Ο ΦΠΑ είναι $372-300=72$ €.\par
\apantisi Αξία $300$ € και ΦΠΑ $72$ €.\par
\elegxos $\frac{24}{100}\cdot300=72$ και $\frac{24}{124}\cdot372=72$.

\paradeigma{Η αρχική τιμή}
\tixero Μετά από έκπτωση $25\%$ ένα σακίδιο κοστίζει $45$ €.\par
\tithelo Την αρχική τιμή.\par
\possindeo $0{,}75x=45$, $x=45:0{,}75=60$.\par
\apantisi $60$ €.\par
\elegxos Το $25\%$ του $60$ είναι $15$ και $60-15=45$.

\paradeigma{Η ποσοστιαία μεταβολή}
\tixero Η τιμή ενός σακιδίου ανεβαίνει από $40$ σε $52$ € και μετά πέφτει στα $39$ €.\par
\tithelo Το ποσοστό κάθε μεταβολής.\par
\possindeo Πρώτη: $\frac{52-40}{40}=\frac{12}{40}=0{,}3=30\%$ αύξηση. Δεύτερη: $\frac{52-39}{52}=\frac{13}{52}=0{,}25=25\%$ μείωση.\par
\apantisi $+30\%$ και $-25\%$.\par
\elegxos $40\cdot1{,}3=52$ και $52\cdot0{,}75=39$.

\paradeigma{Ο τόκος με κεφαλαιοποίηση}
\tixero Κατάθεση $2\,000$ € με επιτόκιο $3\%$ τον χρόνο, για $2$ έτη, με κεφαλαιοποίηση ανά έτος.\par
\tithelo Ο συνολικός τόκος και το τελικό ποσό.\par
\possindeo Έτος $1$: τόκος $0{,}03\cdot2\,000=60$, κεφάλαιο $2\,060$. Έτος $2$: τόκος $0{,}03\cdot2\,060=61{,}8$. Σύνολο τόκων $60+61{,}8=121{,}8$.\par
\apantisi Τόκοι $121{,}80$ €, τελικό ποσό $2\,121{,}80$ €.\par
\elegxos Χωρίς κεφαλαιοποίηση οι τόκοι θα ήταν $2\cdot60=120$ €, λίγο λιγότεροι. Σωστό.

\tmima{Ασκήσεις}

\zoni{Το κατάλαβα}{εκπτώσεις, αυξήσεις, μεταβολές, τόκοι}

\ex Εφάρμοσε έκπτωση $20\%$ σε: α) $45$ € \quad β) $12$ € \quad γ) $130$ €.

\ex Εφάρμοσε αύξηση $15\%$ σε: α) $40$ € \quad β) $200$ € \quad γ) $18$ €.

\ex Βρες την ποσοστιαία μεταβολή: α) από $50$ σε $60$ \quad β) από $80$ σε $60$ \quad γ) από $25$ σε $35$ \quad δ) από $120$ σε $90$.

\ex Βρες την αρχική τιμή: α) μετά από αύξηση $20\%$ έγινε $72$ \quad β) μετά από έκπτωση $40\%$ έγινε $45$ \quad γ) το $15\%$ της είναι $12$.

\ex Η τιμή ενός προϊόντος με ΦΠΑ $24\%$ είναι $248$ €. Βρες τον ΦΠΑ και την αξία χωρίς ΦΠΑ. Κάνε το ίδιο για τελική τιμή $62$ €.

\ex Βρες τον τόκο: α) $1\,500$ € με επιτόκιο $4\%$ για ένα έτος \quad β) το ίδιο κεφάλαιο για $6$ μήνες.

\zoni{Το χρησιμοποιώ}{στα μαγαζιά και στην τράπεζα}

\ex \textbf{Οι εκπτώσεις του Νίκου.} Τα αθλητικά των $90$ € έχουν έκπτωση $30\%$ και επιπλέον $20\%$. Πόσο κοστίζουν; Πόσο θα κόστιζαν με έκπτωση $50\%$ και πόσο είναι η διαφορά;

\ex \textbf{Το μπουφάν και οι μπότες.} Ένα μπουφάν των $120$ € έχει έκπτωση $15\%$ και μια μπότα των $85$ € έχει έκπτωση $20\%$. Πόσα πληρώνει η Αγγελική για τα δύο μαζί;

\ex \textbf{Οι αυξήσεις.} Ένα ποδήλατο των $250$ € ανεβαίνει κατά $10\%$ και τον επόμενο μήνα κατά ακόμη $10\%$. Πόσο κοστίζει; Πόσο ποσοστό αύξηση είναι συνολικά;

\ex \textbf{Η κατάθεση.} Κατάθεση $5\,000$ € με επιτόκιο $2\%$ και κεφαλαιοποίηση ανά έτος. Βρες τον τόκο και το ποσό μετά από $2$ έτη.

\ex \textbf{Το κεφάλαιο.} Τι κεφάλαιο πρέπει να καταθέσω με επιτόκιο $3\%$ για να πάρω $60$ € τόκο σε ένα έτος;

\ex \textbf{Ο λογαριασμός.} Ο λογαριασμός κινητού έχει πάγιο $12$ € και χρέωση $0{,}05$ € ανά μονάδα. Για $800$ μονάδες και ΦΠΑ $24\%$ επί του συνόλου, πόσα πληρώνει η Φανή;

\ex \textbf{Οι δόσεις.} Ένας φορητός υπολογιστής κοστίζει $800$ €. Η Πολυξένη πληρώνει το $25\%$ μετρητά και το υπόλοιπο σε $3$ μηνιαίες δόσεις με τόκο $1\%$ τον μήνα επί του υπολοίπου. Πόσο είναι ο συνολικός τόκος και τι ποσοστό της αρχικής τιμής;

\ex \textbf{Το διάγραμμα.} Η ημέρα του Στέλιου έχει $24$ ώρες: ύπνος $8$, σχολείο $6$, διάβασμα $3$, αθλητισμός $2$, υπόλοιπα $5$. Βρες τα ποσοστά (στο δέκατο) και σχεδίασε ένα ραβδόγραμμα όπου $10$ εκ. αντιστοιχούν στο $100\%$.

\ex \textbf{Η επιστροφή.} Μια τιμή έπεσε κατά $20\%$. Πόσο \% πρέπει να ανέβει για να γυρίσει στην αρχική;

\zoni{Το σκέφτομαι}{ποσοστά που μπερδεύουν}

\ex Ο Βασίλης λέει ότι έκπτωση $30\%$ και μετά $20\%$ κάνουν $50\%$. Πάρε μια τιμή $100$ € και δείξε το λάθος.

\ex Μια τιμή πέφτει κατά $50\%$ και μετά ανεβαίνει κατά $50\%$. Είναι ίδια με την αρχική;

\ex Η προσφορά «$2+1$ δώρο» (πληρώνεις $2$, παίρνεις $3$) τι ποσοστό έκπτωσης είναι;

\ex Σε ένα κυκλικό διάγραμμα φαίνονται τα ποσοστά $28\%$, $35\%$, $12\%$ και ένα τέταρτο ποσοστό. Πόσο είναι το τέταρτο;

\ex Γιατί ο ΦΠΑ μέσα σε μια τιμή με ΦΠΑ $24\%$ δεν είναι το $24\%$ της τελικής τιμής; Δώσε παράδειγμα με $100$ € αξία.

\zoni{Για το σπίτι}{πρόσεξε τις προσφορές}

\ex Βρες μια διαφήμιση με έκπτωση και υπολόγισε την τελική τιμή. Συμφωνεί με αυτήν που αναγράφεται;

\ex Φτιάξε ένα κυκλικό διάγραμμα για το πώς χρησιμοποιείς το χαρτζιλίκι σου.
